\section{Discussion}\label{sec:discussion}

\paragraph{When do model internals help an agent skip its validator?} Across two domains and five models, internal
signals were most useful when the agent's failures were failures to consult information that is in the prompt. In path
planning, a wrong step is almost always a step whose adjacency line the model did not attend to, and both
region-based grounding and the broader internal signals added 5--15 points of skipped validator calls on top of strong
observable baselines. In CSTR fault recovery, the relevant measurements and physics are short and salient; the model
attends to them and errs in the direction or size of its correction. There, the observable state and the proposed
action carry most of the information, and internal signals added little except for the smallest model. For
practitioners, the lesson is to diagnose the agent's failure mode before investing in internal-signal screens:
attention-based grounding pays off for retrieval-heavy tasks with long, structured contexts.

\paragraph{A practical recipe} Three design choices mattered more than the choice of signal. First, accept
thresholds must be chosen with a confidence bound, not a point estimate, or the realised failure rate among unchecked
accepts exceeds the tolerance on new data. Second, a screen trained on first proposals must not be applied to
retries without validation: retry prompts contain failure feedback, which shifts internal signals, and two of five
models' internal probes accepted failing retries; the first-proposal-only rule removed this at no measurable cost in
recovery. Third, the time saved depends on what a skipped call is replaced with. An unchecked accept saves a full
validation, but an unchecked reject triggers a regeneration, so screens save time only when validation is more
expensive than generation. This is exactly the regime that motivates selective validation: high-fidelity or
long-horizon simulators, scenario sets, or multi-unit plants, where validation cost grows while the agent's cost
stays fixed.

\paragraph{Plant models versus model internals} A natural objection is that a screen on plant measurements and
proposed setpoints is a learned surrogate of the plant rather than a statement about the agent. The plant-readings-only
baseline addresses this directly: for three models it skips nearly as many calls as the full observable screen
(the outcome is largely determined by the fault state), whereas for Qwen2.5-3B and Llama-3.2-3B it skips almost
none. More importantly, the internals-only screen uses no plant information at all, yet with the first-proposal-only
rule it avoided 29--65\% of validator calls while executing at most two failing actions unchecked in
2000 closed-loop episodes. The agent's own internal state therefore carries reliable information about when
its proposal can be rejected without simulation, independent of any plant model, even though it rarely suffices to
accept proposals unchecked.

\paragraph{Limitations} The study uses small and medium open-weight models (1.5--7B parameters), one of them
quantised, on a single laptop GPU; larger models may fail differently. The CSTR is one simulated plant with four
fault families, one prompt, and a paused plant during decisions; we make no claim about decision latency under real
time pressure. The validator doubles as the plant simulator in closed loop. The CSTR certification sets were too small
to certify accepts at a 10\% tolerance, as anticipated; larger certification sets are the remedy. The Qwen2.5-3B CSTR
screens come from an exploratory freeze. The probes were trained on first proposals only; training on retries is a
natural extension. Compute figures are specific to our hardware, but the re-costing makes the dependence on
validator cost explicit and transferable.
